\chapter{Port Nyanzaru}
\sessioninfo{Session 5}{Mardi 15 septembre 2026}

L'Effleureuse poursuivit sa route vers le sud. Waterdeep disparut derrière l'horizon, puis la Côte des Épées s'éloigna à son tour jusqu'à ne plus être qu'une ligne incertaine. Bientôt, il ne resta autour du trois-mâts que la mer, le vent et les ordres de Corvis. Le voyage jusqu'à Chult devait durer une vingtaine de jours.

Les quatre aventuriers profitèrent des premières heures pour découvrir le navire. Sous le pont principal se trouvaient la cabine du capitaine, deux cabines pouvant accueillir chacune deux personnes et un espace commun où prendre les repas. Plus bas, la cale abritait des tonneaux, des caisses et les marchandises transportées par l'équipage. Quelques canons complétaient l'ensemble, même si Corvis reconnut qu'ils n'avaient pas servi depuis longtemps.

L'Effleureuse pouvait être manœuvrée par un équipage réduit, mais pas raisonnablement par un homme seul. Corvis voyageait d'ordinaire avec plusieurs marins. Cette fois, il avait été prévenu trop tard : son frère Arlan était parti avec un autre bâtiment et le capitaine avait dû se contenter de quatre passagers totalement inexpérimentés.

Corvis leur parla également de Leilon. Autrefois simple port de pêche, la petite ville avait grandi avec l'activité des Faiseurs de Veuves. La guilde avait rallumé la forge magique, attirant peu à peu voyageurs, commerçants et aventuriers. Certains venaient apporter les matières premières nécessaires à son fonctionnement, d'autres négocier des objets magiques ou solliciter directement les Faiseurs de Veuves. Corvis et Arlan, pêcheurs à l'origine, avaient accompagné cette transformation jusqu'à posséder leur propre compagnie de transport.

\begin{narrateur}
Était-ce une bonne situation, pêcheur devenu commerçant ?

Il n'y avait pas de bonne ou de mauvaise situation. Seulement un métier qui sentait moins le poisson.
\end{narrateur}

Obélik s'interrogea sur Chult, mais le capitaine ne put guère le renseigner. Il n'était jamais descendu aussi loin et préférait habituellement les trajets plus courts le long de la côte. De ce qu'il en savait, l'île était sauvage, ses habitants parfois étranges et les commerces peu nombreux. Il avait toutefois accepté de rendre service et les conduirait jusqu'à Port Nyanzaru avant de repartir vers les îles Solaires.

La conversation finit naturellement par revenir à la pêche. Corvis proposa de leur apprendre cet art subtil, dont la première étape consistait à lancer une ligne puis à attendre. Morrÿs révéla alors qu'il pouvait respirer sous l'eau. Cette aptitude inspira aussitôt à ses compagnons une méthode plus directe : l'attacher à une corde et le faire glisser derrière le navire.

La leçon de pêche se transforma en projet de ski nautique.

Personne ne passa à l'exécution.

\scenebreak

À la tombée de la nuit, les cinq occupants de l'Effleureuse organisèrent des quarts afin qu'une personne reste constamment éveillée. Corvis leur imposa une seule règle : quiconque montait sur le pont dans l'obscurité devait attacher sa ligne de vie. La mer paraissait calme, mais il n'était pas question de perdre un membre de l'équipage pour une vague imprévue.

Destos prit le premier quart. La lune était absente et l'humain ne distinguait presque rien au-delà du navire. Il demeura près du poste de navigation, écoutant le reflux des vagues, le vent dans les voiles et les craquements réguliers de la coque. Rien ne troubla sa surveillance.

Dans les cabines, le sommeil des autres fut moins paisible.

Une ombre s'étendit sur leurs rêves.

Morrÿs revit une part de son passé se déformer sous une obscurité oppressante. Obélik fut ramené vers un événement qu'il aurait préféré laisser derrière lui, mais tout y paraissait désormais plus sombre, comme si ses souvenirs s'enfonçaient lentement dans la pénombre. Kael, lui, vit la caravane de son enfance s'engager dans un tunnel interminable, dépourvu de toute sortie.

\scenebreak

Destos réveilla Morrÿs pour le quart suivant. Le nain émergea de son cauchemar couvert de sueur. Destos, qui n'avait encore rien vu d'anormal, ne trouva aucune explication religieuse à son malaise. Il laissa Morrÿs rejoindre le pont et alla prendre sa place dans la cabine.

Morrÿs s'attacha près du mât, puis employa sa magie pour prendre l'apparence d'un véritable capitaine. Ainsi déguisé, il surveilla les flots avec tout le sérieux que méritait cette nouvelle fonction. Sa vision perçait les ténèbres bien mieux que celle de Destos, mais il ne découvrit rien d'autre que le mouvement des vagues et le balancement des cordages.

\illustration[0.55\textwidth]{2026-09-15/001b7893-dd2c-4e13-b04d-9cc650462dfc.png}
{\textit{Morrÿs profite de son quart pour prendre l'apparence d'un véritable capitaine.}}
{morrys_capitaine}

Pendant ce temps, le sommeil rattrapa le clerc. L'ombre entra à son tour dans ses rêves. Il revit le temple qu'il avait quitté et la raison de son départ, mais le lieu semblait se flétrir autour de lui, lentement gagné par une fin sinistre.

Kael remplaça Morrÿs, s'attacha à son tour et parcourut prudemment le pont. Il prit soin de tourner régulièrement dans l'autre sens pour ne pas enrouler sa ligne autour du mât. L'orc ne vit rien d'inquiétant et finit par aller réveiller Obélik, dont les ronflements contredisaient quelque peu les affirmations selon lesquelles il n'arrivait pas à dormir.

Le goliath monta avec une torche. La nuit était toujours paisible et les vagues particulièrement régulières. Leur rythme finit pourtant par avoir raison de son estomac. Obélik se pencha par-dessus le bastingage et nourrit les poissons avant d'aller réveiller Corvis pour le dernier quart.

\begin{narrateur}
Au moins, la prochaine séance de pêche demanderait peut-être moins d'attente.
\end{narrateur}

\scenebreak

Le lendemain matin, une odeur de pain grillé accueillit les aventuriers dans le mess. Corvis les félicita pour cette première nuit : personne n'était tombé à la mer et ils n'avaient presque pas vomi. À bord de l'Effleureuse, cela semblait suffire à constituer un début encourageant.

Les jours suivants s'écoulèrent au rythme des manœuvres, des repas et des quarts. Les quatre compagnons apprirent à répondre plus rapidement aux ordres du capitaine. La côte disparut entièrement et, pendant de longues journées, l'horizon ne leur offrit plus que des vagues à perte de vue.

Vers le quinzième jour, la température commença à grimper. L'air se chargea d'une humidité toujours plus lourde et l'eau prit des teintes cristallines et turquoise. Obélik crut soudain apercevoir un dragon des mers et appela Morrÿs pour le lui montrer. Kael observa à son tour et distingua bientôt ce qui nageait réellement dans le sillage du navire : un groupe de tortues marines. Elles jouèrent dans les vagues pendant près d'une demi-heure avant de reprendre leur route.

La question de savoir si les tortues pouvaient se pêcher fut rapidement posée. Quelqu'un rappela cependant qu'une divinité tortue existait peut-être quelque part et qu'il serait imprudent de se la mettre à dos. Corvis déconseilla lui aussi l'expérience, puis admit en avoir déjà mangé après s'être retrouvé échoué sur une île. Le souvenir qu'il en gardait mêlait surtout l'iode, le sel et un peu de sable. Son frère n'était peut-être pas le meilleur des cuisiniers dans ce genre de situation.

Après vingt jours de traversée, la terre reparut enfin.

\scenebreak

Chult se dessina d'abord comme une masse verte au milieu de l'océan. Une immense montagne dominait le centre de l'île, dont les pentes disparaissaient sous une végétation épaisse. En approchant par le nord, les voyageurs distinguèrent un point plus animé sur la côte : un port protégé par une longue digue qui formait une baie artificielle.

L'Effleureuse passa entre les défenses de Port Nyanzaru. Une lourde chaîne pouvait fermer l'accès à la baie. Sur la droite, un phare dominait l'entrée du port ; sur la gauche, un fort et ses nombreux canons surveillaient les navires. Corvis guida son trois-mâts à travers l'ouverture et rejoignit les quais.

La chaleur les enveloppa dès leur entrée dans le port. L'air était épais, saturé d'humidité et chargé d'odeurs inconnues : bois vert, végétation dense et épices venues des étals voisins. Des pirogues, des barges marchandes et des bâtiments de toutes tailles s'entassaient le long des quais.

Pourtant, l'animation semblait retenue.

Des drapeaux funèbres pendaient de mât en mât.

Une procession avançait dans la rue principale. Des porteurs pieds nus, vêtus de blanc et couverts de cendre, soutenaient un cercueil de bois brut. Derrière eux, des bannières noires étaient tendues à intervalles réguliers et des fleurs blanches reposaient devant les portes. Les étals demeuraient fermés. Dans cette cité éclatante de couleurs, aucun cri ne s'élevait, seulement le pas mesuré du cortège au milieu d'un silence pesant.

\illustration[0.95\textwidth]{2026-09-15/arc-001.png}
{\textit{Le groupe découvre Port Nyanzaru au passage d'une procession funèbre.}}
{arrivee_port_nyanzaru}

Corvis leur indiqua les quartiers de la ville. Les auberges les plus convenables se trouvaient, selon lui, près du bazar, à l'est. Wakanga O'Tamu, leur commanditaire, résidait dans le quartier marchand, à l'ouest. Le capitaine ne pouvait pas les accompagner : il devait trouver un équipage avant de reprendre la mer.

\illustration[0.82\textwidth]{2026-09-15/nyanzaru-map-player-version.png}
{\textit{La carte de Port Nyanzaru découverte par les aventuriers.}}
{carte_port_nyanzaru}

Avant son départ, Kael lui demanda s'il existait des coutumes particulières à respecter. Corvis ne connaissait pas suffisamment la ville pour donner des règles précises. Il leur rappela seulement que Port Nyanzaru constituait le dernier véritable refuge avant la jungle et leur conseilla de ne pas se mettre ses habitants à dos.

Puis il retourna à ses affaires.

Un domestique s'approcha presque aussitôt du groupe. Un éventail plié à la main, il les observa avec une politesse empressée avant de leur demander à voix basse s'ils étaient les voyageurs envoyés par la guilde. Obélik répondit par l'affirmative avant que ses compagnons aient le temps de rappeler qu'ils n'étaient pas censés se présenter comme membres des Faiseurs de Veuves.

L'homme parut satisfait. Wakanga les attendait et sa demeure n'était pas loin. Les aventuriers acceptèrent de le suivre, mais la soudaineté de cette rencontre les incita à rester sur leurs gardes.

\scenebreak

Le chemin traversa d'abord les abords d'un vaste marché où les étals semblaient prospères, puis s'éleva vers un quartier d'une tout autre richesse. Les villas se succédaient derrière une profusion de tentures colorées et de plantes tropicales. Leurs murs étaient couverts de fresques représentant d'immenses lézards, tantôt bipèdes, tantôt quadrupèdes, dont certains découvraient des rangées de dents impressionnantes.

La demeure de Wakanga rivalisait de luxe avec le quartier général des Faiseurs de Veuves. Le domestique conduisit le groupe dans une antichambre, annonça les visiteurs venus chercher le construct, puis disparut. Quelques instants plus tard, le prince marchand vint lui-même les accueillir.

Wakanga O'Tamu confirma qu'il les attendait pour la recherche du Gardien Bouclier. Avant d'en parler, cependant, il devait leur présenter une affaire plus urgente. Sa position à Port Nyanzaru lui imposait une certaine prudence : il attendait d'eux autant de discrétion que d'efficacité et ne souhaitait pas voir sa réputation compromise.

\illustration[0.45\textwidth]{2026-09-15/Wakanga_Otamu_pose.png}
{\textit{Wakanga O'Tamu, prince marchand de Port Nyanzaru.}}
{wakanga_otamu}

Il les conduisit dans un petit salon où les attendait une elfe très âgée. Une capuche encadrait son visage et un masque de dentelle en dissimulait une grande partie. Seules ses oreilles dépassaient sur les côtés. Malgré la dignité de sa posture, la fatigue et les ravages du temps étaient impossibles à ignorer.

Elle se nommait Syndra.

\illustration[0.45\textwidth]{2026-09-15/Syndra_Sylvane_portrait_-_maudit.png}
{\textit{Syndra, l'amie de longue date de Wakanga, frappée par un mal mystérieux.}}
{syndra}

Wakanga la présenta comme une amie de longue date, liée à lui par une confiance qui dépassait la simple amitié. Il aurait donné ses propres richesses pour la sauver, mais aucun traitement connu ne parvenait à enrayer le mal qui la consumait. Syndra prit alors la parole d'une voix affaiblie.

Des années auparavant, elle avait été aventurière. Une faute d'inattention lui avait coûté la vie au cours d'une embuscade, mais d'autres voyageurs avaient retrouvé son corps et l'avaient ramenée dans un temple. Des prêtres étaient parvenus à la ressusciter, et elle avait poursuivi ses aventures comme si la mort n'avait été qu'une halte sur sa route.

Jusqu'à maintenant.

Depuis quelque temps, son corps dépérissait. Chaque jour lui retirait un peu de ses forces et la rapprochait de nouveau de la mort. Elle n'était pas la seule : toutes les personnes autrefois ramenées à la vie semblaient désormais frappées du même mal. Pire encore, les prêtres ne parvenaient plus à ressusciter les morts, quels que soient leur puissance ou les rites employés.

Syndra avait enquêté aussi longtemps que son état le lui avait permis. La source du phénomène se trouvait quelque part dans les jungles de Chult, probablement loin au sud, mais elle n'avait pas réussi à la localiser. D'autres groupes étaient partis avant eux. Aucun n'était revenu.

Le dernier avait quitté Port Nyanzaru une vingtaine de jours plus tôt.

Le mal était récent. Il progressait vite.

Le temps pressait.

\scenebreak

La mort occupait une place importante dans les croyances locales. Les processions devenaient d'autant plus somptueuses que le défunt avait été puissant ou fortuné, et les prêtres avaient longtemps prolongé certaines existences ou rappelé les morts lorsque leurs moyens le permettaient. Désormais, même eux échouaient. Les cérémonies redoublaient de faste dans l'espoir d'apaiser les divinités, sans que personne sache si celles-ci étaient responsables.

Syndra, elle, soupçonnait moins une malédiction divine qu'un artefact magique. Quelque chose s'était réveillé dans la jungle. Quelque chose qu'il faudrait retrouver et arrêter.

Wakanga leur confia alors la rumeur qui circulait entre les guildes de Port Nyanzaru. Un objet appelé le \textit{Soulmonger} se trouverait quelque part sur l'île et capturerait les âmes des défunts. Il ignorait encore si cette information était exacte. Si l'artefact existait, certains chercheraient certainement à le détruire tandis que d'autres voudraient s'en emparer.

Le prince marchand aurait aimé l'étudier et découvrir s'il était possible d'en neutraliser les effets. Il ne cachait ni son intérêt ni sa confiance dans ses propres talents de mage. Pourtant, il reconnut que certains objets étaient trop dangereux pour être contrôlés. Si aucun autre choix n'existait, le \textit{Soulmonger} devrait être détruit.

Les aventuriers acceptèrent d'aider Syndra. La recherche du Gardien Bouclier restait leur mission initiale, mais lorsqu'Obélik demanda laquelle des deux tâches devait primer, la réponse vint naturellement : une vie importait davantage qu'un objet. Le construct pourrait attendre. Syndra, elle, ne disposait peut-être déjà plus de beaucoup de temps.

Wakanga leur promit toutes les ressources qu'il pourrait mettre à leur disposition. Leur décision lui parut d'autant plus satisfaisante que la demande de son amie était désormais aussi la sienne.

Deux missions les attendaient dans la jungle.

Et aucune ne leur indiquait précisément où chercher.

\scenebreak

Wakanga leur conseilla de ne pas quitter la ville trop rapidement. Chult était immense, dangereuse et encore largement inexplorée. Avant de partir, ils devraient recueillir des informations, acheter le matériel nécessaire et engager un guide capable de les conduire dans la jungle. S'aventurer au hasard au-delà de Port Nyanzaru reviendrait à se perdre avant même d'avoir commencé leurs recherches.

La cité était gouvernée par sept princes marchands. Chacun disposait d'une voix et contrôlait une partie du commerce ou de l'administration. Wakanga gérait les objets magiques et l'information. Ekene-Afa s'occupait notamment des armes, des boucliers et du matériel de voyage ; Ifan Talro'a des animaux et de leur dressage ; Jessamine des plantes et des ordres d'assassinat ; Jobal des guides et des espions ; Kwayothé de la nourriture, des boissons, des parfums et des répulsifs ; Zhanthi des gemmes, des vêtements et des armures.

Le fonctionnement de la justice locale retint particulièrement leur attention. L'esclavage était interdit et les crimes ordinaires sévèrement punis. Un affront personnel pouvait être réglé à l'amiable ou porté devant les autorités. Il était même possible, sous certaines conditions et contre paiement, d'obtenir légalement un ordre d'assassinat.

Les autres condamnés risquaient la course contre la mort. Le criminel était placé dans une enceinte située hors de la ville et devait atteindre la sortie tandis que des créatures choisies selon la gravité de sa faute étaient lâchées à ses trousses. Un vélociraptor pouvait suffire pour les délits les moins graves. Les pires pouvaient mériter beaucoup plus imposant.

Quiconque survivait voyait sa dette envers la cité effacée.

La foule, elle, pouvait parier sur le résultat.

Wakanga leur assura néanmoins que les étrangers ne risquaient pas de provoquer un incident pour une coutume anodine qu'ils ignoraient. Il suffisait de ne pas voler, tuer, blesser ou insulter ceux qui ne leur avaient rien fait. En cas de conflit, mieux valait s'adresser aux princes marchands que tenter de régler seul le problème.

Kael revint aux fresques aperçues dans les rues. Wakanga lui apprit qu'elles représentaient les dinosaures de Chult. Ces animaux appartenaient à la faune locale, mais certains étaient aussi dressés pour participer à des courses organisées dans la cité. Les écuries victorieuses et leurs montures devenaient de véritables célébrités, capables d'enrichir ceux qui avaient parié sur elles.

Il existait plusieurs catégories, des hadrosaures et dimétrodons jusqu'aux jeunes tricératops et tyrannosaures. Les épreuves restaient mouvementées, mais les accidents mortels étaient, selon Wakanga, suffisamment rares.

Le dernier ne remontait qu'à deux mois.

Une nouvelle course devait avoir lieu dans deux jours. Ils pourraient y assister, parier ou même tenter de monter eux-mêmes l'un des dinosaures. La proposition éveilla suffisamment de curiosité pour ne pas être immédiatement écartée.

\scenebreak

Wakanga déploya ensuite une carte de Chult. De vastes régions y demeuraient inconnues, même de lui et des explorateurs qu'il avait envoyés. À l'échelle de l'île, Port Nyanzaru ne semblait être qu'une porte minuscule ouverte sur une immensité verte.

\illustration[0.95\textwidth]{2026-09-15/chult-map-player-version.png}
{\textit{La carte de Chult remise au groupe par Wakanga. Les régions masquées restent inexplorées ; chaque hexagone représente environ quinze kilomètres.}}
{carte_chult}

Le journal d'une naine offrait leur première piste concernant le Gardien Bouclier. D'après ses indications, le construct se trouverait dans une région située à l'intérieur des terres. Suivre la rivière depuis Port Nyanzaru pourrait faciliter une partie du voyage, mais l'itinéraire exact restait à déterminer.

Iolka et Nera Vesper auraient pu les renseigner. Malheureusement, Wakanga ne les avait pas vues depuis une vingtaine de jours. Les deux sœurs n'accompagnaient pas le groupe parti pour Syndra : elles accomplissaient une autre mission dont il ignorait la nature. Elles passaient habituellement par sa villa pour donner de leurs nouvelles, car sa demeure constituait un refuge sûr pour les aventuriers avec lesquels il travaillait.

Leur silence était inhabituel, mais pas encore forcément alarmant. Les expéditions pouvaient durer longtemps.

\begin{narrateur}
Pas de nouvelles, bonnes nouvelles.
\end{narrateur}

Wakanga recommanda au groupe la maison de repos de Kaya, une auberge située près du bazar rouge, et promit d'en couvrir les frais. Il leur indiqua également où trouver le matériel et les produits exotiques dont ils auraient besoin. La ville commerçait en épices, fruits, ivoire et bien d'autres marchandises venues de l'île.

Il insista surtout sur les dangers les moins spectaculaires de la jungle. Les plantes toxiques, les créatures hostiles et les tribus dangereuses n'étaient pas les seules menaces. Se perdre, manquer d'eau potable ou succomber à une maladie loin de tout secours pouvait tuer aussi sûrement qu'un prédateur. Et tant que le mal qui frappait Syndra persistait, nul ne pouvait compter sur une résurrection en cas d'échec.

Destos se montra confiant dans ses capacités à repérer les maladies. Il pourrait notamment vérifier si les points d'eau étaient sûrs. Wakanga accueillit cette assurance avec prudence et lui recommanda une nouvelle fois de ne rien sous-estimer.

Syndra rappela alors que, s'ils mouraient, personne n'était certain de pouvoir les ramener.

L'ironie arracha tout de même un rire à la vieille elfe.

\scenebreak

Les quatre compagnons quittèrent finalement la villa et prirent la direction du quartier du marché. Ils avaient un toit, une carte et plusieurs pistes, mais aucune route certaine. Devant eux s'étendait une ville inconnue ; au-delà de ses murs, une jungle dont presque personne ne revenait indemne.

Kael observa les rues autour d'eux. Il ne remarqua rien de précis parmi les passants, les tentures et la végétation qui envahissait les façades.

Pourtant, une sensation persistait.

Quelqu'un les observait.
